\chapter{\rm\bfseries Introduction}
\label{ch:introduction}

\section{Motivation}

Fabric handling automation is a long-standing open problem in industrial
robotics, with applications spanning apparel manufacturing and textile
logistics, laundry processing and folding, hospital bedding management,
assistive robotics for elderly or disabled individuals, and the packaging
of soft goods. Yet a striking fraction of published robot-manipulation
demonstrations sidesteps the hardest step entirely: the object is handed
to the robot's gripper by a human operator, and the demonstrated
autonomy begins only after the grasp already
exists.\footnote{A prominent recent example: a widely circulated
January 2024 video of the Tesla Optimus humanoid folding a shirt was
subsequently acknowledged by its poster to be teleoperated, with the
shirt pre-staged in a fixed setup~\cite{musk2024optimus,novak2024magic}.} For cloth
specifically, acquiring the first grasp---picking a thin, flat fabric
directly off a surface---remains the bottleneck that limits how
autonomous these systems can truly be. Cloth is uniquely difficult to
grasp robustly: it deforms continuously under contact, exhibits no fixed
graspable feature, and has a near-zero bending stiffness that allows it to slip
out of a poorly placed gripper. Conventional rigid-jaw grippers, designed for
sturdy objects with well-defined geometry, perform poorly on cloth without
elaborate tactile sensing or human-tuned approach trajectories.

A class of compliant grippers based on the finray effect---a passive
bending mechanism inspired by fish-fin anatomy---has been proposed as a
solution. Finray fingers passively wrap around the contact surface, increasing
contact area and gripping force without requiring active force control. This
compliance, however, comes at the cost of force-modelling difficulty: the
finger geometry deforms under load, the contact point shifts, and tilt-induced
asymmetries arise that have no analogue in the rigid-jaw case.

The central question of this thesis is: \emph{to what extent does the finray's
compliance trade single-axis force precision for multi-axis operational
robustness?} Specifically, how does the finray's contact force respond to
deviations in the two design degrees of freedom most likely to be uncontrolled
in real deployment: press depth and gripper tilt angle?

\section{Research Question and Contributions}

This thesis answers the following research questions through a controlled
experimental benchmark on a Franka FR3 collaborative robot --- a benchmark
that treats the gripper itself as the experimental variable, holding the
robot, control policy, cloth, and environment identical across all
conditions:
\begin{enumerate}
  \item How does stationary contact force vary with press depth for each of
        six candidate gripper variants?
  \item For each variant, how much can the gripper be tilted (about each
        orthogonal axis) before the grasp becomes unreliable?
  \item Is the dependence of grasp quality on tilt itself depth-dependent
        --- and is the response direction-dependent (handed) about the
        finray's geometrically asymmetric axis?
  \item What is the compliant finger's mechanical response under
        compression, and is that response shared across the fin-ray
        design family?
\end{enumerate}

\noindent The contributions of this work are:
\begin{itemize}
  \item A reproducible benchmark protocol for cloth-grasp force characterisation
        on a Franka FR3, with full source code published alongside this thesis.
  \item A two-dimensional depth--tilt experimental matrix sampling six
        gripper variants: two rigid parallel-jaw configurations and four
        fin-ray prints spanning blade thickness, seat angle, and print
        instance.
  \item Empirical evidence of a depth-dependent tilt robustness transition
        in finray grippers, with quantitative characterisation of the
        transition depth.
  \item A documented calibration procedure for per-variant tool-centre-point
        (TCP) offset using differential force touch-off against a rigid
        reference.
  \item A geometric verification framework that disentangles control-frame
        positioning error from physical gripper compliance under load.
\end{itemize}

\paragraph{Scope.} The benchmark deliberately isolates the \emph{contact}
phase of cloth grasping---the interval between first gripper--cloth contact
and a completed lift---and characterises it over a continuous depth--tilt
parameter space. Downstream task-level evaluation (e.g.\ folding a cloth
diagonally from a corner grasp, dropping a centre-grasped cloth onto a
curved object, or dragging a cloth between positions under the same
human-defined policy) was part of the original project plan but is
descoped from this thesis; it is discussed as future work in
Chapter~\ref{ch:conclusion}. This narrowing is intentional: it is
precisely because the contact phase is usually bundled inside end-to-end
task pipelines that its mechanics have not previously been mapped as a
function of gripper design.

\section{Thesis Organisation}

Chapter~\ref{ch:background} reviews the literature on cloth grasping,
compliant grippers, and force-controlled approach strategies.
Chapter~\ref{ch:methods} describes the experimental hardware, the kinematic
calibration procedure, and the data-acquisition pipeline.
Chapter~\ref{ch:experiments} presents the depth-sweep characterisation (Experiment~1)
across all five gripper variants at zero tilt.
Chapter~\ref{ch:experiments} presents the tilt-sweep characterisation (Experiment~2)
at boundary and safe depths.
Chapter~\ref{ch:discussion} synthesises the findings into a model of
depth-dependent tilt robustness.
Chapter~\ref{ch:conclusion} concludes and identifies directions for future work.
